% CS388 Final Project report: skeleton.
% Template: NeurIPS 2026 (neurips_2026.sty from the official ZIP, in this folder).
% Option "preprint": shows the author name, no line numbers, footer reads "Preprint."
% (The default anonymizes the author and adds line numbers; "final" needs a track
% option and would print "40th Conference on NeurIPS" in the footer, which a
% course project should not claim.)
% Page budget (course rule, overrides the template's 9): 6-8 pages excluding references.
% Target lengths are in comments.

\documentclass{article}
\usepackage[preprint]{neurips_2026}

\usepackage[utf8]{inputenc}
\usepackage[T1]{fontenc}
\usepackage{hyperref}
\hypersetup{colorlinks=true, linkcolor=black, citecolor=blue!60!black, urlcolor=blue!60!black}
\usepackage{url}
\usepackage{booktabs}
\usepackage{amsfonts,amsmath}
\usepackage{nicefrac}
\usepackage{microtype}
\usepackage{xcolor}
\usepackage{graphicx}
\graphicspath{{../figures/}}

\title{When Can a Medical QA Model's Confidence Be Trusted?\\
Comparing Token-Probability, Verbalized, and Sampling-Based Calibration on MedMCQA}

\author{Troy Ramsay \\ The University of Texas at Austin \\ \texttt{tramsay@utexas.edu}}

\begin{document}
\maketitle

\begin{center}
  \textbf{Video:} \url{https://youtu.be/REPLACE-ME}
\end{center}

\begin{abstract}
% ~150-200 words. Problem, track, what was compared, 3 headline numbers, takeaway.
TODO
\end{abstract}

\section{Introduction}
% ~0.75 page.
% - Motivation: confidence is only useful if it means what it says; medical QA raises the stakes.
% - State the track explicitly: "This project falls under the Uncertainty and Calibration track."
% - Research questions (Q1 methods, Q2 fixes, Q3 scale, Q4 where it breaks, Q5 selective answering).
% - Contributions (bullet list); note the H2 result that contradicts prior work.
TODO

\section{Related Work}
% ~0.75 page. Guo et al. 2017; Desai & Durrett 2020; Kadavath et al. 2022;
% Xiong et al. 2024; Tian et al. 2023; Zheng et al. 2024; Wang et al. 2023;
% Kamath et al. 2020; Nixon et al. 2019; Pal et al. 2022 (MedMCQA).
TODO

\section{Method}
% ~1 page.
% - Definition of confidence: probability that the selected answer letter is correct.
% - Three estimators: token probability (A-D renormalized), verbalized (0-100), sampling agreement (k=10, T=0.7, anchored on greedy answer).
% - Letter-bias correction: mean over 4 cyclic option rotations.
% - Post-hoc calibration: temperature scaling (token prob.), Platt scaling (verbalized, sampling), fit on calibration split.
% - Metrics: ECE (15 equal-width bins; primary), AUROC, Brier, selective AUC, accuracy at coverage; paired bootstrap.
TODO

\section{Experimental Setup}
% ~1 page.
% - MedMCQA validation (4,183) -> exclusions (244 cross-referencing, 28 image-dependent, 2 Unknown) -> 3,909.
% - Stratified split by subject: 1,000 calibration / 2,909 test, seed 388 (reproduced exactly on a fresh machine).
% - Models: Qwen3.5-4B and Qwen3.5-27B, bf16, vLLM, thinking mode off, zero-shot.
% - Compute: single A40 (4B) and A100 80GB (27B).
% - Hypotheses H1-H7 (could be a compact table).
TODO

\section{Results}
% ~2 pages: main table (both models), reliability figure, coverage figure, letter-bias figure, model comparison.
TODO

\section{Analysis}
% ~1.5 pages: by training stage (dental), error categories, method disagreement, hypothesis-by-hypothesis.
TODO

\section{Limitations}
% ~0.5 page: contamination; MedMCQA label noise and digitization artifacts; pattern-based image filter;
% one model family; multiple choice and zero-shot only; ECE upward bias; associations, not causes.
TODO

\section{Conclusion}
% ~0.25 page.
TODO

% \nocite{*} lists every refs.bib entry until the text cites them; remove it
% once all seven papers are cited in the text.
\nocite{*}
\bibliographystyle{plainnat}
\bibliography{refs}

\end{document}
